\documentclass[12pt]{article}
\usepackage[utf8]{inputenc}
\usepackage[UTF8]{ctex}
\usepackage{biblatex}
\usepackage{amssymb}
\usepackage{latexsym}
\usepackage{amsmath}
\usepackage{cases}
\usepackage{geometry}
\usepackage{graphicx}
\usepackage{float}
\usepackage{listings}
\usepackage{enumerate}
\usepackage{color}
\usepackage{minted}


\lstset{ 
  backgroundcolor=\color{white},   % choose the background color; you must add \usepackage{color} or \usepackage{xcolor}
  basicstyle=\ttfamily,            % the size of the fonts that are used for the code
  breakatwhitespace=false,         % sets if automatic breaks should only happen at whitespace
  breaklines=true,                 % sets automatic line breaking
  captionpos=b,                    % sets the caption-position to bottom
  commentstyle=\ttfamily\color{mygreen},    
                                   % comment style
  deletekeywords={},               % if you want to delete keywords from the given language
  escapeinside={},                 % if you want to add LaTeX within your code
  extendedchars=true,              % lets you use non-ASCII characters; for 8-bits encodings only, does not work with UTF-8
  frame=single,                    % adds a frame around the code
  keepspaces=true,                 % keeps spaces in text, useful for keeping indentation of code (possibly needs columns=flexible)
  keywordstyle=\color{blue},       % keyword style
  language=C++,                    % the language of the code
  morekeywords={library,use},                 % if you want to add more keywords to the set
  numbers=left,                    % where to put the line-numbers; possible values are (none, left, right)
  numbersep=5pt,                   % how far the line-numbers are from the code
  numberstyle=\tiny\color{mygray}, % the style that is used for the line-numbers
  rulecolor=\color{black},         % if not set, the frame-color may be changed on line-breaks within not-black text (e.g. comments (green here))
  showspaces=false,                % show spaces everywhere adding particular underscores; it overrides 'showstringspaces'
  showstringspaces=false,          % underline spaces within strings only
  showtabs=false,                  % show tabs within strings adding particular underscores
  stepnumber=1,                    % the step between two line-numbers. If it's 1, each line will be numbered
  stringstyle=\color{mymauve},     % string literal style
  tabsize=2,                       % sets default tabsize to 2 spaces
  title=\lstname                   % show the filename of files included with \lstinputlisting; also try caption instead of title
}
\lstdefinelanguage{VHDL}{
  morekeywords={
    library,use,all,entity,is,port,in,out,end,architecture,of,
    begin,and
  },
  morecomment=[l]--
}
\usepackage{xcolor}
\colorlet{keyword}{blue!100!black!80}
\colorlet{comment}{green!90!black!90}
\lstdefinestyle{vhdl}{
  language     = VHDL,
  basicstyle   = \ttfamily,
  keywordstyle = \color{keyword}\bfseries,
  commentstyle = \color{comment}
}

\addbibresource{bib.bib}
\setlength{\parindent}{0em}
\bibliography{bib}
\geometry{a4paper,scale=0.8}
\linespread{1.5}
\graphicspath{{Images/}}

\begin{document}

\begin{titlepage}
\begin{center}
\linespread{1.2}\huge {\bfseries Fudan University}\\[1cm]
\linespread{1}
\includegraphics[width=10cm]{fudan-name.pdf}\\[3cm]
\linespread{1.9}\huge {\bfseries 可编程器件与硬件描述语言}\\
\linespread{1.9}\LARGE {\bfseries 作业报告}\\[2.5cm]

{\Large 姓名：名字}\\[0.3cm]
\Large 学号：123456789\\[4cm] 
\large 2021年11月1日
\end{center}
\end{titlepage}



\section{优先编码器}
设计一个BCD的优先编码器电路，输入为10个开关的状态，要求输出开
关对应的编码。输出编码用4位表示，第一个开关为0时，输出为0000时，第二个开关为0时，输出为0001时，...... 第10个开关为0时，输出为1001。第10个开关的优先级最高。当没有按键按下时，输出信号E为1。有按键按下时，输出信号E为0。
%题干
\subsection{顶层说明}
\begin{figure}[H]
\centering
\includegraphics[width=8cm]{Q1TOP.png} %这里换成你的图片
\caption{元件示意图}
\end{figure}
Input[9:0]为十个开关，其中Input(9)的优先级最高，Output[3:0]为四位输出，E表示是否有按键被按下。
\begin{figure}[H]
\centering
\includegraphics[width=15cm]{Q1RTL.png}
\caption{RTL图}
\end{figure}
\subsection{每个宏单元如何实现（对所写代码的说明）}
代码如下：
 \begin{minted}[mathescape,
               linenos,
               numbersep=5pt,
               gobble=0,
               frame=single,
               fontfamily=courier,
               framesep=4mm]{vhdl}
---------------------------------------------------
-- 10-bit Priority encoder
-- by Zhang Ke, 2021/10/30 9:26PM
---------------------------------------------------

library ieee;
use ieee.std_logic_1164.all;

---------------------------------------------------

entity BCD_Encoder10 is
  port (
    Input  : in std_logic_vector(9 downto 0);
    Output : out std_logic_vector(3 downto 0);
    E      : out std_logic
  );
end BCD_Encoder10;

---------------------------------------------------

architecture rtl of BCD_Encoder10 is

begin
  process (Input)
  begin
    if (Input(9) = '0') then
      Output <= "1001";
      E      <= '0';
    elsif (Input(8) = '0') then
      Output <= "1000";
      E      <= '0';
    elsif (Input(7) = '0') then
      Output <= "0111";
      E      <= '0';
    elsif (Input(6) = '0') then
      Output <= "0110";
      E      <= '0';
    elsif (Input(5) = '0') then
      Output <= "0101";
      E      <= '0';
    elsif (Input(4) = '0') then
      Output <= "0100";
      E      <= '0';
    elsif (Input(3) = '0') then
      Output <= "0011";
      E      <= '0';
    elsif (Input(2) = '0') then
      Output <= "0010";
      E      <= '0';
    elsif (Input(1) = '0') then
      Output <= "0001";
      E      <= '0';
    elsif (Input(0) = '0') then
      Output <= "0000";
      E      <= '0';
    else
      Output <= "XXXX";
      E      <= '1';
    end if;
  end process;

end rtl;
  \end{minted}
以上代码中，将元件输入Input定义为10位的标准逻辑向量类型，输出E和Output分别定义为标准逻辑位和4位的标准逻辑向量类型。具体architecture实现中使用IF语句进行状态判断并且输出。\\
\subsection{时序仿真结果}
编写Test Bench文件如下：
 \begin{minted}[mathescape,
               linenos,
               numbersep=5pt,
               gobble=0,
               frame=single,
               fontfamily=courier,
               framesep=4mm]{vhdl}
---------------------------------------------------
-- Test Bench for 10-bit Priority encoder
-- by Zhang Ke, 2021/10/30 9:57PM
---------------------------------------------------

library ieee;
use ieee.std_logic_1164.all;

---------------------------------------------------

entity Q1_BCD_Encoder_tb is
end Q1_BCD_Encoder_tb;

---------------------------------------------------

architecture tbbehv of Q1_BCD_Encoder_tb is

  component BCD_Encoder10
    port (
      Input  : in std_logic_vector(9 downto 0);
      Output : out std_logic_vector(3 downto 0);
      E      : out std_logic
    );
  end component;

  signal Input  : std_logic_vector(9 downto 0) := "0000000000";
  signal Output : std_logic_vector(3 downto 0);
  signal E      : std_logic;
begin

  U_BCD   : BCD_Encoder10 port map(Input, Output, E);
  inputpc : process
  begin
    Input <= "0000000000";  wait for 20 ns; Input <= "1000000000";  wait for 20 ns;
    Input <= "1100000000";  wait for 20 ns; Input <= "1110000000";  wait for 20 ns;
    Input <= "1111000000";  wait for 20 ns; Input <= "1111100000";  wait for 20 ns;
    Input <= "1111110000";  wait for 20 ns; Input <= "1111111000";  wait for 20 ns;
    Input <= "1111111100";  wait for 20 ns; Input <= "1111111110";  wait for 20 ns;
    Input <= "1111111111";  wait for 20 ns; Input <= "1001100000";  wait for 20 ns;
    Input <= "1110100000";  wait for 20 ns; Input <= "1111100100";  wait for 20 ns;
    Input <= "1111011100";
    wait; 
  end process;

end architecture;
  \end{minted}
使用上述Test Bench文件进行仿真，波形图如下：
\begin{figure}[H]
\centering
\includegraphics[width=17.5cm]{Q1BL.png}
\caption{BCD优先编码器时序仿真波形}
\end{figure}


\end{document}